%% sections/06-discussion.tex

\section{Discussion}
\label{sec:discussion}

\subsection{The Amplification Mechanism}

The amplification mechanism we propose is: resource exhaustion, when legible and
emotionally architectured, \textit{forces character decisions to be visible}. In a
no-pressure session, a PC can resolve a situation through tactical default (attack
until the encounter ends). Under resource pressure, tactical defaults become costly:
every spell slot spent is permanently gone; every HP taken persists. The party must
decide whether the tactical use is worth the narrative cost.

When the resource is also emotionally architectured (it represents people, not just
power), the visibility of character decisions is higher: conserving healing supplies
is not just resource-management, it is a decision about who the party is willing to
let go without help. In Campaign 4, Sava (the healer PC) managing slots across 7
sessions as a unified resource rather than session-by-session is an example: her
slot decisions at session 3 were explicitly shaped by her grief-paragraph's pattern
(not spending prematurely, because early spending had cost Elara in her backstory).
The resource mechanic made this grief-paragraph-driven behavior visible as a character
decision rather than letting it dissolve into session-by-session optimization.

\subsection{Why the Cohesion Die Amplifies More Strongly}

The Cohesion Die tracks people rather than materials. Three properties distinguish this:

\begin{enumerate}
  \item \textbf{Named loss:} Each Cohesion Die degradation represents named unit members.
    The party knows who is gone. Material depletion (healing supplies at d4) does not
    have names.

  \item \textbf{Irreversibility:} Unit members who are gone do not return. Healing supplies
    at d4 could theoretically be resupplied; the 4 unit members lost between sessions 3
    and 7 cannot be. Irreversibility increases the weight of each Cohesion Die decision.

  \item \textbf{Active preservation:} The party in Campaign 6 consistently made decisions
    that prioritized keeping unit members alive over tactical advantage. In Sessions 4--6,
    at least one PC per session made a sub-optimal tactical choice (measured by expected
    damage output) to preserve a unit member. This preservation behavior \textit{is itself
    character-driven decision-making}: it reflects the party's commitment to the people
    the Cohesion Die represents, not tactical optimization.
\end{enumerate}

The Supply Die does not produce the same preservation behavior: the party in Campaign 4
used healing supplies when they were needed, not in a pattern of active conservation.
The party treated the Supply Die as a constraint; they treated the Cohesion Die as
a responsibility.

\subsection{Design Implications}

For designers creating resource mechanics for long-form narrative campaigns:

\begin{enumerate}
  \item \textbf{Make the resource legible:} A degrading die that represents a named
    quantity (d12 $\to$ d8 $\to$ d6 $\to$ d4 $\to$ exhausted) is more effective at
    producing character-driven engagement than abstract HP numbers alone. The party
    must be able to see their cumulative depletion at each session.

  \item \textbf{Architect the emotional connection:} Resources that represent people
    (or relationships, or commitments) produce stronger amplification than resources
    that represent materials. Design the resource so that its depletion is felt as
    a loss, not just as a constraint.

  \item \textbf{Allow preservation as character expression:} Design the resource so
    that actively preserving it is a valid (if sub-optimal tactically) character choice.
    If preservation is never rational, the party will not preserve; they will simply
    expend the resource when needed.

  \item \textbf{Do not resupply too early:} In Campaign 4, the Supply Die not recovering
    between sessions (combined with no long rests) was the mechanism that made depletion
    visible. If the party had been able to resupply mid-campaign, the cumulative
    pressure would have reset and the amplification effect would not have accumulated.
\end{enumerate}

\subsection{Limitations}

The primary limitation is the confound between resource mechanic and campaign archetype.
Campaign 4 is a siege scenario; Campaign 6 is a military unit. Both feature sustained-pressure
narratives that may independently activate character-driven decision-making. We cannot
rule out that the CA amplification is driven by archetype (siege/duty produces high CA
in sessions 5--7 regardless of resource mechanics) rather than resource mechanic.

A controlled experiment would require two campaigns of the same archetype, one with
and one without a degrading-resource mechanic, with all other spine components held
constant. The Marathon corpus does not contain this comparison. Campaign 4 and 5
provide a partial control (professional-code archetype, Supply Die vs.\ HP-only),
but the archetypes are still different enough (siege vs.\ professional-crime) that
the comparison is imperfect.

We also note that the amplification effect was observed only in long-form campaigns
(7 sessions). Whether the same mechanic produces amplification in 3--4 session
campaigns (where cumulative depletion cannot reach the same level) is unknown.
